\documentclass[11pt]{article}
% ===========================================================================
%  Shared preamble for all MPM2D worksheets — input right after \documentclass.
%  (Files beginning with "_" are skipped by mpm2d-worksheet-publish.mjs.)
% ===========================================================================
\usepackage[letterpaper,top=2.2cm,bottom=1.9cm,left=1.6cm,right=1.6cm,headheight=28pt,footskip=24pt]{geometry}
\usepackage{amsmath,amssymb}
\usepackage{cancel}
\usepackage[T1]{fontenc}
\usepackage[utf8]{inputenc}
\usepackage{lmodern}
\usepackage{xcolor}
\usepackage[most]{tcolorbox}
\usepackage{enumitem}
\usepackage{multicol}
\usepackage{fancyhdr}
\usepackage{tikz}
\usetikzlibrary{calc}
\usepackage{pgfplots}
\pgfplotsset{compat=1.18}

\definecolor{exblue}{HTML}{4A90E2}
\definecolor{exbg}{HTML}{E6F3FF}
\definecolor{qorange}{HTML}{E69138}
\definecolor{qbg}{HTML}{FFF7CC}
\definecolor{keygray}{HTML}{475569}
\definecolor{keybg}{HTML}{F1F5F9}
\definecolor{iagreen}{HTML}{14653B}
\definecolor{titleblue}{HTML}{1D4ED8}
% Rotating palette so each worked example gets its own colour.
\definecolor{ex0}{HTML}{2563A0}\definecolor{ex1}{HTML}{3B7D3B}\definecolor{ex2}{HTML}{8A5A00}
\definecolor{ex3}{HTML}{A3327A}\definecolor{ex4}{HTML}{6D28A3}\definecolor{ex5}{HTML}{B08900}
\definecolor{ex6}{HTML}{0E7490}\definecolor{ex7}{HTML}{9A3412}\definecolor{ex8}{HTML}{15803D}

\pagestyle{fancy}
\fancyhf{}
\fancyhead[L]{\footnotesize NAME: \underline{\hspace{3.4cm}}}
\fancyhead[C]{\textbf{Integration Academy}\\[-2pt]{\scriptsize Math Simplified \textperiodcentered{} Success Amplified}}
\fancyhead[R]{\footnotesize MPM2D}
\fancyfoot[L]{\scriptsize dr.merie@IntegrationAcademy.ca}
\fancyfoot[C]{\scriptsize\thepage}
\fancyfoot[R]{\scriptsize IntegrationAcademy.ca}
\renewcommand{\headrulewidth}{1.2pt}
\renewcommand{\footrulewidth}{0.4pt}
\setlength{\parindent}{0pt}
\setlength{\parskip}{4pt}

% Examples are NOT breakable, so each example + its solution stays on one page.
\newtcolorbox{example}[2]{colframe=#1,colback=#1!7,coltitle=white,colbacktitle=#1,fonttitle=\bfseries,title={#2},arc=3pt,boxrule=1pt}
\newtcolorbox{qbox}[1]{colframe=qorange,colback=qbg,coltitle=white,colbacktitle=qorange,fonttitle=\bfseries,title={#1},arc=3pt,breakable,boxrule=1pt}
\newtcolorbox{keybox}{colframe=keygray,colback=keybg,coltitle=white,colbacktitle=keygray,fonttitle=\bfseries,title={$\checkmark$\ Answer Key},arc=3pt,breakable,boxrule=1pt}
\newtcolorbox{recapbox}{colframe=keygray,colback=keybg,coltitle=white,colbacktitle=keygray,fonttitle=\bfseries,title={Question Recap},arc=3pt,breakable,boxrule=1pt}
\newtcolorbox{ideabox}{colframe=titleblue,colback=titleblue!6,coltitle=white,colbacktitle=titleblue,fonttitle=\bfseries,title={The Big Ideas},arc=3pt,breakable,boxrule=1.2pt}
\newtcolorbox{learnbox}{colframe=iagreen,colback=iagreen!5,coltitle=white,colbacktitle=iagreen,fonttitle=\bfseries,title={The Big Picture},arc=3pt,breakable,boxrule=1.2pt}
\newcommand{\lh}[1]{\par\smallskip{\bfseries\color{iagreen}#1}\par\nobreak\smallskip}

\newcommand{\soln}{\par\smallskip\textit{\textbf{Solution.}}\ }
\newcommand{\workspace}[1]{\par\vspace{3pt}\fcolorbox{black!30}{white}{\begin{minipage}[t][#1][t]{\dimexpr\linewidth-2\fboxsep\relax}\ \end{minipage}}\par}
\newcommand{\sech}[1]{\par\vspace{8pt}{\Large\bfseries\color{iagreen}#1}\par\nobreak\vspace{2pt}{\color{black!20}\hrule height1pt}\vspace{6pt}}

% A compact coordinate plot sized for an example box: \eplot{xmin}{xmax}{ymin}{ymax}{body}
\newcommand{\eplot}[5]{\begin{center}\begin{tikzpicture}
\begin{axis}[width=6.8cm,height=4.4cm,axis lines=middle,xlabel={\tiny$x$},ylabel={\tiny$y$},
grid=both,grid style={gray!16},xmin=#1,xmax=#2,ymin=#3,ymax=#4,
xtick distance=2,ytick distance=2,tick label style={font=\tiny},axis line style={-},samples=120]
#5
\end{axis}\end{tikzpicture}\end{center}}

% A sine-curve plot (degrees on x, 0..360): \splot{ymin}{ymax}{body}
\newcommand{\splot}[3]{\begin{center}\begin{tikzpicture}
\begin{axis}[width=8.6cm,height=4.6cm,axis lines=middle,xlabel={\tiny$x\,(^\circ)$},ylabel={\tiny$y$},
grid=both,grid style={gray!16},xmin=0,xmax=360,ymin=#1,ymax=#2,
xtick distance=90,ytick distance=1,tick label style={font=\tiny},axis line style={-},samples=180]
#3
\end{axis}\end{tikzpicture}\end{center}}

% A small coordinate plot: \plot{xmin}{xmax}{ymin}{ymax}{body}
\newcommand{\plot}[5]{\begin{center}\begin{tikzpicture}
\begin{axis}[width=8.4cm,height=6cm,axis lines=middle,xlabel={\scriptsize$x$},ylabel={\scriptsize$y$},
grid=both,grid style={gray!16},xmin=#1,xmax=#2,ymin=#3,ymax=#4,
xtick distance=2,ytick distance=2,tick label style={font=\tiny},axis line style={-}]
#5
\end{axis}\end{tikzpicture}\end{center}}

% Right triangle (angle theta at bottom-left, right angle at bottom-right):
%   \rtri{drawBase}{drawHeight}{oppLabel}{adjLabel}{hypLabel}
\newcommand{\rtri}[5]{\begin{center}\begin{tikzpicture}[scale=0.85]
\coordinate (A) at (0,0);\coordinate (B) at (#1,0);\coordinate (C) at (#1,#2);
\draw[thick,exblue] (A)--(B)--(C)--cycle;
\draw[black] ($(B)+(-0.32,0)$)--($(B)+(-0.32,0.32)$)--($(B)+(0,0.32)$);
\node[below] at ($(A)!0.5!(B)$){#4};
\node[right] at ($(B)!0.5!(C)$){#3};
\node[above left] at ($(A)!0.5!(C)$){#5};
\node at (0.7,0.22){$\theta$};
\end{tikzpicture}\end{center}}

% General (oblique) triangle with vertex labels and opposite-side labels:
%   \gtri{Alabel}{Blabel}{Clabel}{aSide}{bSide}{cSide}
\newcommand{\gtri}[6]{\begin{center}\begin{tikzpicture}[scale=0.85]
\coordinate (A) at (0,0);\coordinate (B) at (5,0);\coordinate (C) at (1.6,3);
\draw[thick,exblue] (A)--(B)--(C)--cycle;
\node[below left] at (A){#1};\node[below right] at (B){#2};\node[above] at (C){#3};
\node[below] at ($(A)!0.5!(B)$){#6};
\node[right] at ($(B)!0.5!(C)$){#4};
\node[left] at ($(A)!0.5!(C)$){#5};
\end{tikzpicture}\end{center}}

\fancyhead[R]{\footnotesize MCV4U \textperiodcentered{} 6.1}
\begin{document}

\begin{center}
{\LARGE\bfseries\color{titleblue}Worksheet 6.1 --- Introduction to Vectors}\\[2pt]
{\small Grade 12 Calculus and Vectors (MCV4U) \textperiodcentered{} Unit 6: Geometry and Algebra of Vectors}
\end{center}

\begin{ideabox}
A vector has magnitude and direction (velocity, force); a scalar has magnitude only (speed, mass). Equal vectors share magnitude and direction; $-\vec v$ reverses direction.
\begin{itemize}[leftmargin=*,itemsep=2pt]
  \item Vector = magnitude + direction; scalar = magnitude only.
  \item Equal vectors: same magnitude and direction (position-free).
  \item Directions via bearings ($N30^\circ E$, $030^\circ$).
\end{itemize}
\end{ideabox}

\begin{learnbox}
\lh{Magnitude and direction}A \textbf{vector} has both size and direction, unlike a scalar. Displacement, velocity, and force are vector quantities.
\lh{Representing vectors}Draw a vector as a directed arrow; two vectors are equal when they share magnitude and direction, wherever they are placed.
\lh{Special vectors}The \textbf{zero vector} has no length, and the opposite of a vector reverses its direction.
\end{learnbox}

\sech{Worked Examples}

\begin{example}{ex0}{Example 1 --- Vector or scalar}
Classify velocity, speed, force, temperature, displacement.\soln
\textbf{Step 1 (ask whether direction matters for each quantity):} velocity, force, and displacement all depend on which way they point; speed and temperature do not.\\[3pt]
\textbf{Conclusion:} vectors: velocity, force, displacement. Scalars: speed, temperature.
\end{example}

\begin{example}{ex1}{Example 2 --- Magnitude \& direction}
A plane flies 500 km at $N60^\circ E$.\soln
\textbf{Step 1 (identify the magnitude):} the distance travelled is $500$ km.\\[3pt]
\textbf{Step 2 (identify the direction):} the bearing $N60^\circ E$ means $60^\circ$ measured from north, toward east.\\[3pt]
\textbf{Conclusion:} magnitude $500$ km; direction $N60^\circ E$.
\end{example}

\begin{example}{ex2}{Example 3 --- Equal vectors}
When are two displacement vectors equal?\soln
\textbf{Step 1 (recall what defines a vector):} magnitude and direction, not position.\\[3pt]
\textbf{Conclusion:} two vectors are equal exactly when they share magnitude and direction --- where they're drawn doesn't matter.
\end{example}

\begin{example}{ex3}{Example 4 --- Opposite vector}
If $\vec v$ is 8 m/s east, describe $-\vec v$.\soln
\textbf{Step 1 (recall what $-\vec v$ means):} same magnitude, reversed direction.\\[3pt]
\textbf{Step 2 (apply it):} magnitude stays $8$ m/s; east reverses to west.\\[3pt]
\textbf{Conclusion:} $-\vec v$ is 8 m/s west.
\end{example}

\begin{example}{ex4}{Example 5 --- Bearing}
Express "$40^\circ$ east of north" as a true bearing.\soln
\textbf{Step 1 (recall true bearing format):} measured clockwise from north, as a 3-digit angle.\\[3pt]
\textbf{Step 2 (convert):} "$40^\circ$ east of north" is already $40^\circ$ clockwise from north.\\[3pt]
\textbf{Conclusion:} $040^\circ$.
\end{example}

\begin{example}{ex5}{Example 6 --- Vector or scalar}
Is acceleration a vector or scalar?\soln
\textbf{Step 1 (ask if direction matters):} acceleration has a direction (the way the velocity is changing).\\[3pt]
\textbf{Conclusion:} vector.
\end{example}

\begin{example}{ex6}{Example 7 --- Vector or scalar}
Is mass a vector or scalar?\soln
\textbf{Step 1 (ask if direction matters):} mass is just an amount, with no direction.\\[3pt]
\textbf{Conclusion:} scalar.
\end{example}

\begin{example}{ex7}{Example 8 --- Opposite}
Describe $-\vec v$ if $\vec v$ is 10 N north.\soln
\textbf{Step 1 (keep the magnitude, reverse the direction):} $10$ N stays the same; north reverses to south.\\[3pt]
\textbf{Conclusion:} $-\vec v$ is 10 N south.
\end{example}

\begin{example}{ex8}{Example 9 --- Bearing}
Write $N75^\circ E$ as a true bearing.\soln
\textbf{Step 1 (convert from N-E form to true bearing):} $N75^\circ E$ means $75^\circ$ clockwise from north.\\[3pt]
\textbf{Conclusion:} $075^\circ$.
\end{example}

\clearpage
\sech{Your Turn --- Show Your Work}
For each question, show all steps clearly in the space provided.

\begin{qbox}{Question 1}
Is force a vector or a scalar?\workspace{2.4cm}
\end{qbox}

\begin{qbox}{Question 2}
Is temperature a vector or a scalar?\workspace{2.4cm}
\end{qbox}

\begin{qbox}{Question 3}
Describe $-\vec v$ if $\vec v$ is 6 m/s west.\workspace{2.4cm}
\end{qbox}

\begin{qbox}{Question 4}
Two vectors have equal magnitude, opposite direction. Equal?\workspace{2.4cm}
\end{qbox}

\begin{qbox}{Question 5}
Write $N20^\circ E$ as a true bearing.\workspace{2.4cm}
\end{qbox}

\begin{qbox}{Question 6}
Is displacement a vector or a scalar?\workspace{2.4cm}
\end{qbox}

\begin{qbox}{Question 7}
Is speed a vector or a scalar?\workspace{2.4cm}
\end{qbox}

\begin{qbox}{Question 8}
Describe $-\vec v$ if $\vec v$ is 12 N east.\workspace{2.4cm}
\end{qbox}

\begin{qbox}{Question 9}
What two things define a vector?\workspace{2.4cm}
\end{qbox}

\begin{qbox}{Question 10}
Write $N50^\circ E$ as a true bearing.\workspace{2.4cm}
\end{qbox}

\begin{qbox}{Question 11}
Is energy a vector or a scalar?\workspace{2.4cm}
\end{qbox}

\begin{qbox}{Question 12}
Does the position of a vector affect whether two are equal?\workspace{2.4cm}
\end{qbox}

\begin{qbox}{Question 13 --- Challenge}
A ship sails 30 km at a bearing of $120^\circ$. State its magnitude and describe the bearing in N–E form.\workspace{2.6cm}
\end{qbox}

\clearpage
\sech{Question Recap \& Answer Key}

\begin{recapbox}
\begin{multicols}{2}
\small
\begin{enumerate}[leftmargin=*,itemsep=3pt]
  \item Is force a vector or a scalar?
  \item Is temperature a vector or a scalar?
  \item Describe $-\vec v$ if $\vec v$ is 6 m/s west.
  \item Two vectors have equal magnitude, opposite direction. Equal?
  \item Write $N20^\circ E$ as a true bearing.
  \item Is displacement a vector or a scalar?
  \item Is speed a vector or a scalar?
  \item Describe $-\vec v$ if $\vec v$ is 12 N east.
  \item What two things define a vector?
  \item Write $N50^\circ E$ as a true bearing.
  \item Is energy a vector or a scalar?
  \item Does the position of a vector affect whether two are equal?
  \item A ship sails 30 km at a bearing of $120^\circ$. State its magnitude and describe the bearing in N–E form.
\end{enumerate}
\end{multicols}
\end{recapbox}

\begin{keybox}
\begin{multicols}{2}
\small
\begin{enumerate}[leftmargin=*,itemsep=3pt]
  \item vector
  \item scalar
  \item 6 m/s east
  \item no (opposites)
  \item $020^\circ$
  \item vector
  \item scalar
  \item 12 N west
  \item magnitude and direction
  \item $050^\circ$
  \item scalar
  \item no
  \item 30 km; $120^\circ=S60^\circ E$
\end{enumerate}
\end{multicols}
\end{keybox}

\end{document}
